\part{User space}

\chapter{The TSS and the way back in}

\begin{goals}
  \item Why a ring 3 program cannot run on its own stack inside the kernel
  \item What the TSS is and why only one field of it matters
  \item How the GDT descriptor for the TSS is filled in from C
\end{goals}

\section{The problem}

A user process runs in ring 3 with its own stack, somewhere around \hexa{BFFFF000}. Now the timer
fires. The processor must run kernel code --- but on which stack?

Not the user's. It might be nearly full, it might be corrupt, it might be deliberately pointed at
kernel memory by a hostile program. If the kernel pushed its data there, a user program could read
the kernel's saved state, or make the kernel fault by unmapping its own stack.

So the processor \hi{switches stacks automatically} on a privilege change. And it needs to be
told, in advance, which stack to switch to. That is what the \term{TSS} (Task State Segment) is for.

\begin{analogy}
The user stack is the visitor's own bag. The kernel will not put its confidential paperwork in a
visitor's bag --- the visitor could read it, or walk off with it. So at the door there is a
dedicated locker, and its number is written on a card the receptionist keeps. The TSS is that card;
\cod{esp0} is the locker number.
\end{analogy}

\section{One useful field}

The TSS is a large structure --- 104 bytes --- designed in the 80286 era for hardware task
switching, a feature nobody uses any more because it is slower than doing it in software.

Modern systems, including this one, keep a TSS only because the processor insists on reading two
fields from it:

\begin{center}
\begin{tabularx}{0.95\textwidth}{@{}l X@{}}
\toprule
\textbf{Field} & \textbf{Purpose} \\
\midrule
\cod{ss0} & The stack \emph{segment} to use in ring 0. Always \hexa{10} here \\
\cod{esp0} & \textbf{The stack pointer to use in ring 0.} Updated on every switch to a user process \\
\bottomrule
\end{tabularx}
\end{center}

Everything else in the structure is zeroed and ignored.

\filetag{lytlnybl/tasks/procman.c}
\begin{lstlisting}[style=c]
void init_tss(void) {
    uintptr_t base = (uintptr_t)&tss;
    uint32_t limit = sizeof(tss_t) - 1;

    /* fill in the GDT descriptor that was left blank in assembly */
    gdt_tss[0] = limit & 0xFF;
    gdt_tss[1] = (limit >> 8) & 0xFF;
    gdt_tss[2] = base & 0xFF;
    gdt_tss[3] = (base >> 8) & 0xFF;
    gdt_tss[4] = (base >> 16) & 0xFF;
    gdt_tss[5] = 0x89;              /* present, ring 0, 32-bit available TSS */
    gdt_tss[6] = (limit >> 16) & 0x0F;
    gdt_tss[7] = (limit >> 24) & 0xFF;

    memset(&tss, 0, sizeof(tss_t));
    tss.ss0 = 0x10;                 /* kernel data segment */
    tss.esp0 = 0;                   /* set properly on each context switch */
    tss.iomap_base = sizeof(tss_t); /* no I/O permission bitmap */

    load_tss();
}
\end{lstlisting}

This is the descriptor from Chapter 16 being filled in by hand, byte by byte, at runtime --- because
the address of a C variable is not known at assembly time. Access byte \hexa{89} means present, DPL
0, system segment, type 9 = ``available 32-bit TSS''.

\begin{warn}{\cod{iomap\_base}}
Setting it to \cod{sizeof(tss\_t)} means ``the I/O permission bitmap starts past the end of this
structure'', which the processor reads as ``there is no bitmap, deny all port access from ring 3''.
Leave it at zero and the processor would interpret the start of the TSS itself as a permission
bitmap, granting user programs direct access to arbitrary hardware ports. A subtle and genuinely
dangerous mistake.
\end{warn}

\begin{lstlisting}[style=asm]
load_tss:
    mov ax, 0x28        ; the TSS selector
    ltr ax              ; load task register
    ret
\end{lstlisting}

\section{What actually happens on entry}

\begin{center}
\begin{tikzpicture}[font=\sffamily\scriptsize,
  b/.style={draw=grayborder,rounded corners=3pt,minimum width=3.4cm,minimum height=1.5cm,align=center}]

  \node[b,fill=greenlight,draw=green] (u) at (0,2.4) {\textbf{Ring 3}\\{\tiny ESP = 0xBFFFEF00}\\{\tiny user stack}};
  \node[b,fill=amberlight,draw=amber] (t) at (5.6,2.4) {\textbf{TSS}\\{\tiny esp0 = 0xC0004000}\\{\tiny ss0 = 0x10}};
  \node[b,fill=redlight,draw=red] (k) at (2.8,-0.2) {\textbf{Ring 0}\\{\tiny ESP = 0xC0004000}\\{\tiny kernel stack of this process}};

  \draw[-{Stealth[length=2.5mm]},amber,thick] (u) -- node[above,font=\tiny,text=gray]{int 0x80} (t);
  \draw[-{Stealth[length=2.5mm]},red,thick] (t) -- node[right,font=\tiny,text=gray,align=left]{the CPU reads esp0\\and switches} (k);
  \draw[-{Stealth[length=2.5mm]},green,thick,dashed] (k) -- node[left,font=\tiny,text=gray]{iret} (u);
\end{tikzpicture}
\end{center}

Once on the kernel stack, the processor pushes the user's \reg{SS} and \reg{ESP} there, then
\reg{EFLAGS}, \reg{CS} and \reg{EIP}. That is why \cod{registers\_t} has those two extra fields at
the end, and why they only exist when the ring changed.

% ============================================================
\chapter{Entering ring 3}

\begin{goals}
  \item How a process is placed in ring 3 without a special instruction
  \item How a program's address space is built
  \item How a binary is loaded from disk into memory
\end{goals}

\section{There is no ``enter user mode'' instruction}

x86 has no instruction that says ``drop to ring 3''. Privilege is lowered only as a side effect of
returning from an interrupt: \cod{iret} pops \reg{CS}, and if the popped \reg{CS} has RPL 3, the
processor switches to ring 3, pops \reg{SS} and \reg{ESP} too, and resumes there.

So the trick is to fabricate a context that \emph{looks like} it came from ring 3, and then let the
normal return path do the work:

\filetag{lytlnybl/tasks/procman.c}
\begin{lstlisting}[style=c]
void create_uprocess(process_t* new_process) {
    new_process->regs.esp = USER_STACK_TOP;      /* 0xBFFFF000 */
    new_process->regs.cs = 0x18 | 3;             /* user code, RPL 3 */
    new_process->regs.ds = 0x20 | 3;             /* user data, RPL 3 */
    new_process->regs.ss = 0x20 | 3;
    new_process->regs.eflags = 0x202;            /* interrupts enabled */
    /* ... address space set-up ... */
}
\end{lstlisting}

When the scheduler picks this process, \cod{load\_context} copies those values into the
\cod{registers\_t} on the kernel stack, and \cod{iret} does the rest. \hi{Ring 3 is entered by
returning to it, from an interrupt that never happened.}

\begin{concept}{Why \reg{EFLAGS} must have \reg{IF} set}
The user process must run with interrupts enabled, otherwise the timer can never preempt it and a
single \cod{while(1);} would freeze the machine permanently --- with no way to recover, because
ring 3 cannot execute \cod{sti} either.

Note this is not a favour the kernel does the program: it is the kernel protecting \emph{itself}.
\end{concept}

\section{Building the address space}

\begin{lstlisting}[style=c]
new_process->page_directory = (page_directory_t*)alloc_frame();
memset((page_directory_t*)new_process->page_directory, 0, 4096);

for (uint32_t i = 0; i < 1024; i++) {
    (*new_process->page_directory)[i] = (*kernel_directory)[i];
}

uintptr_t ustack_frame = (uintptr_t)alloc_frame();
map_page(new_process->page_directory, USER_STACK_TOP - 4096, ustack_frame,
         PAGE_PRESENT | PAGE_WRITABLE | PAGE_USER);

map_page(new_process->page_directory, USER_HEAP_START,
         (uintptr_t)alloc_frame(),
         PAGE_PRESENT | PAGE_WRITABLE | PAGE_USER);

new_process->user_heap_end = USER_HEAP_START + 4096;
new_process->heap_pages_allocated = 1;
new_process->ustack = (void*)(USER_STACK_TOP - 4096);
\end{lstlisting}

One page of stack (4 KB) and one page of heap. Note the stack page is mapped at
\cod{USER\_STACK\_TOP - 4096}, because \cod{USER\_STACK\_TOP} is the address \emph{above} the last
usable byte --- the stack grows down from there into the page below.

\begin{danger}{The stack does not grow}
Exactly one page. If a user program recurses deeply or declares a large local array, \reg{ESP}
drops below \hexa{BFFFE000} and hits an unmapped page. The page fault handler prints a message and
the process is killed.

Real systems handle this: a fault just below the stack is recognised as legitimate growth and a new
page is mapped automatically. Adding that to \cod{page\_fault\_handler} is one of the most
instructive extensions you can make to this project --- it is perhaps fifteen lines, and it turns a
crash into a feature.
\end{danger}

\section{Loading a program from disk}

\filetag{lytlnybl/tasks/loader.c}
\begin{lstlisting}[style=c]
process_t* load_program(const char* path) {
  int32_t fd = fs_open(path);
  uint32_t inode_num = fs_fd_to_inode(fd);

  fs_inode_t file_inode;
  if (!fs_read_inode(inode_num, &file_inode))
    return NULL;

  size_t file_size = file_inode.size;

  uint8_t buffer[file_size];                    /* variable length array */
  if (fs_read(fd, buffer, file_size) == -1)
    return NULL;

  process_t* user_proc = create_process((void*)USER_CODE_BASE, PROCESS_USER);

  for (size_t i = 0; i < ((file_size + 4095) / 4096); ++i) {
    void* code_frame = alloc_frame();

    size_t offset = i * 4096;
    size_t bytes = file_size - offset;
    if (bytes > 4096) bytes = 4096;
    memcpy(code_frame, buffer + offset, bytes);

    map_page(user_proc->page_directory, USER_CODE_BASE + (4096 * i),
             (uintptr_t)code_frame,
             PAGE_PRESENT | PAGE_WRITABLE | PAGE_USER);
  }

  fs_close(fd);
  return user_proc;
}
\end{lstlisting}

The logic is: read the whole file into a buffer, then copy it page by page into freshly allocated
frames, mapping each one at consecutive virtual addresses starting at \hexa{400000}.

\begin{concept}{Why the copy is needed at all}
\cod{memcpy(code\_frame, ...)} writes to a \emph{physical} address, which works because the kernel
identity-maps the low 4 MB. The mapping into the process's address space happens afterwards.

The kernel cannot write directly to \hexa{400000} in the process's space, because that address means
something else entirely in the kernel's own address space. This two-step dance --- write to the
frame physically, then map it virtually --- is a pattern you will meet in every kernel.
\end{concept}

\begin{danger}{Three real problems in this function}
\textbf{1. \cod{uint8\_t buffer[file\_size]} is a variable-length array on a 16 KB kernel stack.} A
program larger than about 14 KB overflows it, and Chapter 4 explained what that means. The fix is to
copy page by page directly from the file, never materialising the whole thing.

\textbf{2. \cod{fs\_open} is not checked.} If the path does not exist it returns $-1$, and the code
proceeds to call \cod{fs\_fd\_to\_inode(-1)}.

\textbf{3. Code pages are mapped writable.} \cod{PAGE\_WRITABLE} on a code page means a program can
modify its own instructions. Dropping that flag would be a one-word improvement in safety.
\end{danger}

\begin{danger}{And the consequence of problem 2, in the shell}
When the shell runs \cod{run /nonexistent}, \cod{load\_program} returns \cod{NULL}, and the system
call handler does:
\begin{lstlisting}[style=c]
process_t* child = load_program(path);
child->parent_pid = current_process->pid;   /* NULL dereference */
\end{lstlisting}
\cod{parent\_pid} is at offset 4 in the structure, so this writes to address \hexa{4} --- which
\emph{is} mapped, because the low 4 MB is identity-mapped. No fault occurs. Instead it silently
corrupts the interrupt vector table area, then reads a garbage pid, sets the shell to
\cod{BLOCKED} waiting for a child that will never exist, and switches away.

\hi{The observable symptom is that the shell hangs forever with no error message.} Verified on
a running system: after \cod{run /bin/usertest} (a typo), the prompt never returns and further
typing echoes but does nothing.

The fix is three lines:
\begin{lstlisting}[style=c]
process_t* child = load_program(path);
if (!child) { regs->eax = -1; break; }
\end{lstlisting}
\end{danger}

% ============================================================
\begin{recipe}{loading a program from the kernel}{ch33-loader}
Calls \cod{load\_program} twice: once with a path that does not exist, to show it returns
\cod{NULL} as it should, and once with a real one. It then prints the new process's page directory,
entry point, stack pointer and code segment before the scheduler has run it even once.

\cod{CS} comes back as \hexa{1B}, which is \hexa{18}\cod{|3}. Nothing else has to happen: the timer
interrupt reaches the scheduler, the scheduler picks the process, and \cod{iret} drops the CPU into
ring 3. You can watch the program's output appear a moment later.
\end{recipe}

% ============================================================
\chapter{System calls}

\begin{goals}
  \item Follow a system call from the C library to the kernel and back
  \item Understand how arguments and return values travel
  \item Understand why validating user pointers is the most important missing piece
\end{goals}

\section{The user side}

\filetag{lytlnybl/user/libc/syscalls.asm}
\begin{lstlisting}[style=asm]
syscall:
    mov eax, [esp + 4]   ; service number
    mov ebx, [esp + 8]   ; first argument
    mov ecx, [esp + 12]  ; second argument
    mov edx, [esp + 16]  ; third argument

    int 0x80

    ret                  ; EAX already holds the return value
\end{lstlisting}

Ten lines that bridge two worlds. C pushes arguments on the stack; this routine moves them into the
registers the kernel expects, triggers the interrupt, and returns --- with \reg{EAX} untouched since
the kernel put the result there, which is exactly where C expects a return value.

\filetag{lytlnybl/user/libc/syscalls.c}
\begin{lstlisting}[style=c]
int write(int fd, void *buff, size_t count) {
  return syscall(SYSCALL_WRITE, fd, (uint32_t)buff, count);
}

int read(int fd, void* buff, size_t count) {
  return syscall(SYSCALL_READ, fd, (uint32_t)buff, count);
}

void* sbrk(intptr_t increment) {
  return (void*)syscall(SYSCALL_SBRK, increment, 0, 0);
}
\end{lstlisting}

\section{The kernel side}

\filetag{lytlnybl/kernel/interrupts.asm}
\begin{lstlisting}[style=asm]
syscall_entry:
    push dword 0x80
    push dword 0

    pusha
    mov ax, ds
    push eax

    push esp
    call syscall_handler
    add esp, 4

    pop eax
    popa
    add esp, 8
    iret
\end{lstlisting}

Structurally identical to the ISR stub. The two initial pushes fill the \cod{interrupt\_number} and
\cod{error\_code} slots so the layout matches \cod{registers\_t}.

\begin{warn}{A small inconsistency}
The ISR macros push the error code first and the interrupt number second, so the number lands at the
lower address. Here the order is reversed: \hexa{80} ends up in \cod{error\_code} and 0 in
\cod{interrupt\_number}.

It is harmless, because \cod{syscall\_handler} reads neither field. But it will confuse anyone
debugging with a generic dump routine, who will see interrupt number 0 --- ``divide by zero'' --- on
every system call. Swapping the two pushes would cost nothing and remove a trap.
\end{warn}

\section{The dispatcher}

\filetag{lytlnybl/kernel/interrupts.c}
\begin{lstlisting}[style=c]
void syscall_handler(registers_t* regs) {
    int fd;
    uint32_t buffer;
    size_t count;

    switch (regs->eax) {
        case SYSCALL_GETPID:
            regs->eax = current_process->pid;
            break;

        case SYSCALL_YIELD:
            context_switch(current_process, get_next_process(), regs);
            break;

        case SYSCALL_SLEEP:
            current_process->wake_tick = timer_get_ticks() + regs->ebx;
            current_process->state = PROCESS_SLEEPING;
            context_switch(current_process, get_next_process(), regs);
            break;
        /* ... */
    }
}
\end{lstlisting}

\cod{regs->eax = current\_process->pid;} deserves a second look. It writes into the saved state on
the kernel stack; \cod{popa} restores it into the real \reg{EAX}; \cod{iret} returns to ring 3; the
user's \cod{syscall} routine executes \cod{ret} with the value already in place. The return value
travels four layers without anyone explicitly returning anything.

\section{\cod{sbrk}: growing the user heap}

\begin{lstlisting}[style=c]
case SYSCALL_SBRK:
    current_process->user_heap_end += regs->ebx;

    while (((current_process->user_heap_end + 4096) - USER_HEAP_START) / 4096
            > current_process->heap_pages_allocated) {
        map_page(current_process->page_directory,
            USER_HEAP_START + (current_process->heap_pages_allocated++ * 4096),
            (uintptr_t)alloc_frame(),
            PAGE_PRESENT | PAGE_USER | PAGE_WRITABLE);
    }

    regs->eax = current_process->user_heap_end;
    break;
\end{lstlisting}

This is the interface Unix has used since the 1970s: ``move the top of my heap by this much and tell
me where it now is''. The kernel deals only in pages; the user-space \cod{malloc} in
\filepath{user/libc/memory.c} turns those pages into byte-sized allocations, using the same
linked-list design as the kernel heap in Chapter 29.

\begin{warn}{\cod{alloc\_frame} can return NULL}
If physical memory is exhausted, \cod{map\_page} maps address 0 as a user page. The program then
gets a heap that silently aliases physical frame 0 --- the interrupt vector table. Checking the
return value and failing the call cleanly is a small fix with a large payoff.
\end{warn}

\section{The validation that is missing}

This is the most important section in the chapter.

\begin{lstlisting}[style=c]
case SYSCALL_WRITE:
    fd = regs->ebx;
    buffer = regs->ecx;
    count = regs->edx;

    if (fd < 1) { regs->eax = -1; break; }
    if (fd > 2) { regs->eax = fs_write(fd, (void*)buffer, count); break; }

    ((char*) buffer)[count] = '\0';        /* <-- here */
    vga_text_write(&terminal, (char*)buffer);

    regs->eax = (int)count;
    break;
\end{lstlisting}

\begin{danger}{Two distinct problems in one line}
\textbf{1. It writes one byte past what the caller asked for.} If a program calls
\cod{write(1, buf, sizeof(buf))} --- which the shell's \cod{cat} command does --- the kernel writes
a zero one byte beyond the end of \cod{buf}, corrupting whatever follows on the user's stack.

\textbf{2. The pointer is never validated.} \cod{buffer} is a number chosen entirely by the user
program. Nothing checks that it points into that process's own memory. A program can pass a kernel
address and the kernel --- running in ring 0, where the \cod{PAGE\_USER} bit does not restrict it
--- will happily write a zero byte there.

That is a privilege escalation primitive. Not a theoretical one: an arbitrary-address single-byte
write into kernel memory is enough to corrupt a function pointer or a page table entry.
\end{danger}

\begin{concept}{What a real kernel does instead}
Every operating system has a function like \cod{copy\_from\_user} / \cod{copy\_to\_user}. It:
\begin{enumerate}
  \item checks that the whole range $[p, p+n)$ lies inside user address space,
  \item checks the pages are present and have the right permissions,
  \item copies into a kernel buffer, so the user cannot change the data after validation,
  \item returns an error rather than faulting if anything is wrong.
\end{enumerate}
Adding even a minimal version of this to lytlnyblOS --- ``reject any pointer at or above
\hexa{C0000000}, and any \cod{count} that would cross out of a mapped page'' --- would close the
entire class of bug. It is the single highest-value change one could make to this codebase.
\end{concept}

\section{The full round trip}

\begin{center}
\begin{tikzpicture}[font=\sffamily\scriptsize,
  r/.style={draw=grayborder,rounded corners=2pt,minimum width=6.6cm,minimum height=0.62cm,align=left,inner xsep=8pt}]
  \node[r,fill=greenlight,draw=green] (a) {\texttt{printf("hi")} in shell.c};
  \node[r,fill=greenlight,draw=green,below=0.2cm of a] (b) {\texttt{putchar} $\rightarrow$ \texttt{write(1, \&c, 1)}};
  \node[r,fill=greenlight,draw=green,below=0.2cm of b] (c) {\texttt{syscall(5, 1, ptr, 1)} in syscalls.c};
  \node[r,fill=amberlight,draw=amber,below=0.2cm of c] (d) {\texttt{syscall:} load EAX/EBX/ECX/EDX, \texttt{int 0x80}};
  \node[r,fill=redlight,draw=red,below=0.2cm of d] (e) {CPU: ring 3 $\rightarrow$ 0, stack from TSS, jump via IDT};
  \node[r,fill=redlight,draw=red,below=0.2cm of e] (f) {\texttt{syscall\_entry:} build registers\_t};
  \node[r,fill=redlight,draw=red,below=0.2cm of f] (g) {\texttt{syscall\_handler:} case 5 $\rightarrow$ \texttt{vga\_text\_write}};
  \node[r,fill=redlight,draw=red,below=0.2cm of g] (h) {\texttt{regs->eax = count}};
  \node[r,fill=amberlight,draw=amber,below=0.2cm of h] (i) {\texttt{popa}, \texttt{iret}: ring 0 $\rightarrow$ 3};
  \node[r,fill=greenlight,draw=green,below=0.2cm of i] (j) {\texttt{ret} in syscall, EAX is the result};
  \foreach \x/\y in {a/b,b/c,c/d,d/e,e/f,f/g,g/h,h/i,i/j} \draw[-{Stealth[length=2mm]},gray] (\x)--(\y);
\end{tikzpicture}
\end{center}

% ============================================================
\begin{recipe}{adding a system call, end to end}{ch34-new-syscall}
The single most useful recipe in the book, because it is the one you will want first. It adds
\cod{SYSCALL\_UPTIME}, which returns the tick count, and the patch script makes exactly four edits
--- the four you always have to make:

\begin{enumerate}
  \item a number in \filepath{include/kernel/interrupts.h},
  \item a \cod{case} in \cod{syscall\_handler} that writes the result into \cod{regs->eax},
  \item the same number plus a prototype in \filepath{include/user/libc/syscalls.h},
  \item the wrapper in \filepath{user/libc/syscalls.c}.
\end{enumerate}

It ships a user program that calls it, sleeps for 100 ticks and calls it again. Run
\cod{run /bin/user\_test} at the shell and the difference printed is exactly 100.
\end{recipe}

% ============================================================
\chapter{A minimal C library}

\begin{goals}
  \item What a user program needs before \cod{main} can run
  \item How \cod{printf} is implemented with variadic arguments
  \item How user-space \cod{malloc} sits on top of \cod{sbrk}
\end{goals}

\section{What the library provides}

\begin{center}
\begin{tabularx}{\textwidth}{@{}l X@{}}
\toprule
\textbf{File} & \textbf{Contents} \\
\midrule
\filepath{user/libc/syscalls.asm} & The \cod{syscall} bridge \\
\filepath{user/libc/syscalls.c} & One wrapper per system call \\
\filepath{user/libc/output.c} & \cod{putchar}, \cod{puts}, \cod{print\_num}, \cod{printf} \\
\filepath{user/libc/memory.c} & \cod{malloc}, \cod{free}, \cod{calloc}, \cod{realloc}, \cod{memcpy}, \cod{memmove}, \cod{memset}, \cod{memcmp} \\
\filepath{user/libc/string.c} & \cod{strlen}, \cod{strcmp}, \cod{strcpy}\ldots \\
\filepath{user/libc/chars.c} & \cod{isalpha}, \cod{isdigit}, \cod{isspace}, \cod{tolower}\ldots \\
\bottomrule
\end{tabularx}
\end{center}

They are compiled into an archive, \cod{libc.a}, and linked into each program:

\filetag{lytlnybl/Makefile}
\begin{lstlisting}[style=mk]
$(USER_LIB): $(USER_LIB_C_OBJECTS) $(USER_LIB_ASM_OBJECTS)
	$(AR) rcs $@ $^

$(SHELL_ELF): $(BUILD_DIR)/user/programs/shell.o $(USER_LIB)
	$(LD) -m elf_i386 -T $(USER_LINKER) $^ -o $@
\end{lstlisting}

An archive is just a bundle of object files; the linker pulls out only the ones actually referenced,
which keeps the binaries small.

\section{\cod{printf} with variadic arguments}

\filetag{lytlnybl/user/libc/output.c}
\begin{lstlisting}[style=c]
void printf(char* format, ...) {
  va_list args;
  va_start(args, format);

  while (*format) {
    if (*format == '%') {
      format++;
      if (*format == 'c') {
        char c = va_arg(args, int);
        putchar(c);
      } else if (*format == 's') {
        char *str = va_arg(args, char*);
        puts(str);
      } else if (*format == 'd') {
        int num = va_arg(args, int);
        print_num(num);
      } else if (*format == '%') {
        putchar('%');
      } else {
        putchar('%');
        putchar(*format);
      }
    } else {
      putchar(*format);
    }
    format++;
  }
  va_end(args);
}
\end{lstlisting}

\begin{concept}{How variadic arguments work}
Because 32-bit x86 pushes arguments on the stack right to left, the arguments after \cod{format} sit
immediately above it in memory, contiguously. \cod{va\_start} points at that spot and \cod{va\_arg}
reads one and advances.

This also explains why \cod{printf} is fundamentally unsafe: the function has no way to know how
many arguments were actually passed. \cod{printf("\%d \%d", 1)} reads a second integer from whatever
happens to be on the stack. Nothing detects it.
\end{concept}

\begin{warn}{\cod{putchar} writes one byte at a time}
\cod{putchar} calls \cod{write(1, \&c, 1)}, which is a full system call --- an interrupt, a
privilege switch, a dispatch --- \emph{per character}. Printing a 40-character line means 40 system
calls.

Real libraries buffer: accumulate into an array and flush on newline or when full. Adding buffering
here would be a large speed-up and a small change. It would also interact with the write-past-the-end
bug from the previous chapter, so fix that first.
\end{warn}

\section{\cod{malloc} on top of \cod{sbrk}}

\filetag{lytlnybl/user/libc/memory.c}
\begin{lstlisting}[style=c]
void *malloc(size_t bytes) {
  if (bytes == 0) return NULL;
  if (!heap_start_head) init_heap();

  heap_header_t* block = find_free_block(bytes);
  if (block) {
    if (block->size != bytes) split_block(block, bytes);
    block->free = false;
    return (void*)((uintptr_t)block + sizeof(heap_header_t));
  } else {
    expand_heap(bytes);
    return malloc(bytes);
  }
}

void expand_heap(size_t requested_size) {
  uintptr_t original_end = heap_end;
  heap_end = (uintptr_t)sbrk(requested_size);
  /* ... append a new free block ... */
}
\end{lstlisting}

This is the same design as the kernel heap --- headers, a singly linked list, split and merge --- but
one layer up: where the kernel heap called \cod{map\_page}, this calls \cod{sbrk}.

\begin{concept}{The same pattern at every level}
Notice the recurring shape. The physical manager hands out 4 KB frames. The virtual manager turns
frames into mapped pages. The kernel heap turns pages into bytes. \cod{sbrk} turns kernel pages into
user pages. User \cod{malloc} turns those into bytes.

Five layers, each doing one job and asking the layer below for coarser units. That is what
``building an operating system'' looks like in practice: not one clever mechanism, but a stack of
simple ones that each hide the mess of the one beneath.
\end{concept}

\begin{recipe}{a user program template}{ch35-libc}
Everything a program in this system can do, in one commented file: \cod{printf} with all its
conversions, \cod{malloc} / \cod{calloc} / \cod{realloc} / \cod{free}, the string functions, and
\cod{get\_pid}, \cod{sleep}, \cod{yield} and \cod{exit}.

Start from it when writing your own. Note it calls its entry point \cod{\_start}, which is what
\filepath{user.ld} actually names --- unlike the shell, for the reason the next section
explains.
\end{recipe}

\section{The entry point problem}

\filetag{lytlnybl/user/user.ld}
\begin{lstlisting}[style=txt]
ENTRY(_start)

SECTIONS
{
    . = 0x00400000;
    .text   : { *(.text) }
    .rodata : { *(.rodata) }
    .data   : { *(.data) }
    .bss    : { *(.bss) }
}
\end{lstlisting}

The kernel always jumps to \hexa{400000} when starting a user process, so \emph{whatever code sits
at that address} is the entry point.

\begin{danger}{The entry point works by accident}
The linker script declares \cod{ENTRY(\_start)}, and \filepath{user/programs/test.c} does define
\cod{\_start}. But \filepath{user/programs/shell.c} defines \cod{start} --- without the
underscore.

Link the shell and the linker warns:
\begin{lstlisting}[style=txt]
ld: warning: cannot find entry symbol _start; defaulting to 00400000
\end{lstlisting}
It then works anyway, for one reason only: the Makefile lists \cod{shell.o} first, so its
\cod{.text} lands first, and \cod{start()} happens to be the first function in the file.

Add a function above \cod{start()}, or reorder the link line, and the kernel will jump into the
wrong function with no diagnostic at all. The fix is trivial --- rename it to \cod{\_start}, or
better, write a tiny \cod{crt0.o} that is always linked first and calls \cod{main} --- but until
then this is a loaded gun.
\end{danger}

\begin{tryit}
Check it for yourself:
\begin{lstlisting}[style=sh]
i686-elf-nm -n build/shell.elf | head -3
\end{lstlisting}
You should see \cod{U \_start} (undefined) and \cod{00400000 T start}. The address is right; the
symbol name is not.
\end{tryit}

\begin{summary}
User space works: programs run unprivileged, in their own address space, and can ask the kernel for
services through a single controlled door. What is left is giving them something to work with ---
persistent storage --- and a program to drive it all.
\end{summary}
